\subsection{معرفی و انگیزه}
این ماژول در چارچوب پژوهش حاضر پیاده‌سازی و ارزیابی شده و فاقد مقاله مستقل جداگانه در ادبیات تحقیق است.

هرس وزن‌ها (\lr{Weight Pruning}) با هدف کاهش تراکم اتصالات عصبی و کمینه‌سازی میزان مصرف حافظه، نیازمند کنترل دقیق خطای بازسازی است. رویکردهای هرس متداول مبتنی بر قدر مطلق وزن‌ها (\lr{Magnitude Pruning}) عموماً بدون توجه به تعادل بین بازسازی تانسورها و منظم‌سازی تنک عمل می‌کنند که این امر در مقیاس‌های تنک‌سازی بالا منجر به تخریب بازنمایی ویژگی‌های لایه‌ها می‌شود.

ماژول \lr{proxsparse} با نام فارسی «هرس‌کننده تنک گرادیان مجاورتی» به عنوان یکی از اجزای موتور فشرده‌سازی سیستم، مسئله تنک‌سازی را در قالب بهینه‌سازی غیرمحدب مقید با عملگرهای مجاورتی (\lr{Proximal Gradient Optimization}) فرمول‌بندی می‌نماید. این ماژول قادر است الگوهای تنک‌سازی نیمه‌ساختاریافته \lr{N:M} (مانند الگوی استاندارد $2:4$ سخت‌افزاری انویدیا)، هرس بلوکی (\lr{Block-Sparsity}) و هرس نامنظم سراسری را تحت منظم‌سازی نرم $L_1$ هدایت کند.
